% !TeX program = xelatex
\documentclass[UTF8, zihao=-4, fontset=mac]{ctexart} 
\setCJKmainfont[BoldFont=Heiti SC, ItalicFont=Kaiti SC]{Songti SC}
\setCJKsansfont{Heiti SC}
\setCJKmonofont{STFangsong}

% =========================================================
% 1. 宏包与页面全局设置
% =========================================================
\usepackage[a4paper, left=3.5cm, right=3.5cm, top=3cm, bottom=3cm]{geometry} 
\usepackage{setspace}
\setstretch{1.6} % 更加严谨的社科行间距

\usepackage{graphicx}
\usepackage{booktabs} 
\usepackage{amsmath}  
\usepackage{amssymb}
\usepackage{fancyhdr} 
\usepackage{titlesec} 
\usepackage{xcolor}   
\usepackage{enumitem}

% 英文核心字体：Times New Roman
\setmainfont{Times New Roman}
\setsansfont{Arial}

% 参考文献格式：GB/T 7714-2015 顺序编码制
\usepackage[numbers, sort&compress]{gbt7714}

% =========================================================
% 3. 页面装饰 (页眉页脚)
% =========================================================
\pagestyle{fancy}
\fancyhf{}
\fancyhead[C]{\zihao{-5}\songti 生成式治理：地方高校应用型转型困境的复杂系统解析}
\fancyfoot[C]{\zihao{-5}\thepage}
\renewcommand{\headrulewidth}{0.4pt}

% =========================================================
% 2. 标题格式优化 (深度参考《教育研究》等顶刊样式)
% =========================================================
\titleformat{\section}{\centering\zihao{3}\heiti}{}{0em}{\chinese{section}、}
\titleformat{\subsection}{\zihao{-3}\heiti}{}{0em}{（\chinese{subsection}）}
\titleformat{\subsubsection}{\zihao{4}\heiti}{}{0em}{\arabic{subsubsection}.}
\titlespacing*{\section}{0pt}{1.5ex plus .5ex minus .2ex}{1.2ex plus .2ex}
\titlespacing*{\subsection}{2em}{1ex plus .5ex minus .2ex}{0.8ex plus .2ex}


\title{\zihao{2}\heiti\bfseries 生成式治理：地方高校应用型转型困境的复杂系统解析}
\author{\zihao{4}\songti （为匿名评审，作者信息暂不公开）}
\date{}



% =========================================================
% 4. 正文开始
% =========================================================
\begin{document}

\maketitle

\renewcommand{\thefootnote}{\fnsymbol{footnote}}
\footnotetext[1]{本文系全国教育科学规划国家一般课题
“地方高校应用型转型困境及治理策略研究”（批准号：B1A250166）的阶段性成果。}
\renewcommand{\thefootnote}{\arabic{footnote}}


\begin{abstract}
\noindent
{\zihao{5}\heiti 【摘要】}\quad {\zihao{5}\kaishu 引导地方高校向应用型转变是建设教育强国、服务区域高质量发展的关键环节，但实践中长期存在“政策高热”与“微观冷感”并存的悖论。本文依托全国教育科学规划课题（B1A250166）形成的10年纵向面板数据，在复杂适应系统与计算教育学的视角下，构建嵌入大语言模型智能体的“高校治理数字孪生系统”。计算实验表明：单一放大学术绩效工分容易触发“古德哈特效应”；应用型成果等效权重超过30\%--40\%时会出现转型的临界相变；中层干部的风险规避导致政策强度在“最后一公里”衰减。文章提出以“生成式治理”为导向，从经验判断转向基于数字孪生的“政策风洞”，为推进教育治理现代化提供一套可计算、可视化、可迭代的决策支持框架。}

\vspace{0.8em}
\noindent
{\zihao{5}\heiti 【关键词】}\quad {\zihao{5}\kaishu 生成式治理；地方高校转型；计算教育学；复杂适应系统；教育治理现代化}
\end{abstract}


\vspace{1cm}

% =========================================================
% 第一章 问题提出 (Problem Statement)
% =========================================================
\section{问题提出：转型深水区的结构性难题}

\subsection{历史维度的审视：从“规模扩张”到“类型分化”的迟滞}

建设教育强国，龙头在高等教育，基点在地方高校。回溯中国高等教育改革四十余年的演进，不难发现，地方本科院校的“身份紧张”并非近期才出现，而是在大众化、精英化与类型化多重逻辑交织中长期累积的结果\cite{zha2011understanding}。1999 年以来的大规模扩招在很大程度上解决了“有学上”的问题，却在一定程度上遮蔽了“上什么样的学、培养什么样的人”等结构性矛盾：一方面，地方高校数量快速增长，成为高等教育体系的重要承载体；另一方面，其办学定位在“类研究型大学”与“类高职院校”之间长期摇摆。

在相当长时间内，以“硕博点数量”“科研项目级别”“SCI 收录论文”等为代表的学术型指标，成为高校竞争的统一参照系。地方高校出于争取资源与合法性的考虑，普遍采取“向上模仿”的策略，把建设“综合性研究型大学”视为理想路径\cite{pan2014transformation, mok2005globalization}。直接后果是不同层次、不同类型高校之间出现了专业结构、人才培养目标和评价标准的“千校一面”，服务地方经济社会发展的功能被边缘化\cite{zha2011understanding}。

2014 年以来，围绕“应用型大学建设”“地方高校转型发展”等议题，国家和地方陆续出台了一系列政策文件，明确提出要推动一批地方本科院校向应用型转变\cite{pan2014transformation, zhong2014applied}。“坚持类型办学、服务地方发展”在宏观政策层面已形成基本共识。但在具体制度设计上，地方高校长期面临“双重制度逻辑”的张力：一方面，行政文件强调“向下扎根”，突出产教融合、区域服务；另一方面，学位授权、学科评估与科研评价体系仍以学术产出为核心指标，客观上引导高校继续“向上生长”，比肩研究型大学。政策目标的“应用型”与评价体系的“学术型”之间存在显著错位，也为后续的转型实践埋下了制度性张力。

课题组基于 B1A250166 项目对某省地方本科高校的调研显示，在“十三五”“十四五”等中长期规划中，超过三分之二的高校在愿景层面提出“建设高水平应用型大学”，但在职称评审、绩效分配等关键制度文件中，论文和纵向课题仍占据主导权重，应用型成果的评价与激励相对有限。这种“目标陈述上的转型”与“制度结构上的迟滞”，构成了本文分析的基本起点。

\subsection{现实数据的悖论：“形转实不转”的剪刀差}

从纵向数据看，地方高校转型实践中呈现出一个值得警惕的“剪刀差”现象。一方面，各级政府对地方高校的支持力度持续加大：财政性经费稳步增长，地方专项资金、产教融合试点、应用型示范项目等相继落地；另一方面，用以反映应用型转型成效的关键指标（如横向经费占比、技术成果转化率、服务地方项目数量等）提升有限，部分年份甚至出现停滞或波动。

为更直观呈现，这里以某省 2014--2023 年地方本科高校相关统计为例，对比“经费投入增长曲线”与“应用型产出指标曲线”。如图 \ref{fig:paradox} 所示，前者整体呈上升趋势，而后者则相对平缓，两者之间形成不断扩大的“资源—绩效剪刀差”。

% --- 图1占位 ---
\begin{figure}[htbp]
    \centering
    \fbox{\parbox{0.8\textwidth}{\centering \vspace{2cm} \textbf{[图1]：地方高校转型过程中的“资源—绩效”剪刀差示意}\\ \small （示意图：财政投入曲线上升明显，应用型产出曲线相对平缓，二者逐渐分离） \vspace{2cm}}}
    \caption{地方高校转型过程中的“资源—绩效”背离现象（2014--2023）}
    \label{fig:paradox}
    \vspace{-0.5em}
    {\footnotesize \kaishu 数据来源：基于全国教育科学规划课题（B1A250166）调研与监测数据整理。}
\end{figure}

在高校内部微观治理层面，类似张力也有具象体现。一方面，学校层面密集出台与应用型转型相关的制度安排：建设产业学院、鼓励教师承担企业委托项目、提高学生实践学分占比等，形成政策上的“高频信号”；另一方面，一线教师在访谈中普遍反映：“时间主要花在论文和纵向项目上，横向课题和企业合作虽然被鼓励，但在职称和绩效中的折算系数较低，风险和不确定性更大。”这表明：在现有制度结构下，教师群体在主观认同上并不反对应用型转型，但在客观激励上缺乏足够的“行动理由”。

从组织行为视角看，这种“上层制度安排热烈、基层行动者反应有限”的状态，常被概括为“形转实不转”“口号热、行动冷”。若简单归因于“执行不力”或“意识不到位”，既不公平，也难以支撑有效的政策改进。更为合理的解释是：在既有制度、文化与认知结构的共同塑形下，地方高校已经形成较为稳定的“学术偏好路径”。在这一路径上，即便改革信号被持续释放，系统仍表现出显著惯性与迟滞。

\subsection{研究问题与分析视角：从线性因果到复杂系统}

围绕上述历史演进与现实悖论，既有研究已从多种理论视角加以解释。资源依赖理论强调组织对外部关键资源的依赖，帮助我们理解地方高校为何倾向于“迎合上级指标偏好”\cite{mok2005globalization}；新制度主义与组织社会学揭示了制度同形性、合法性追求以及“形式结构的神话化与仪式化”等机制，解释了“规划文件高度一致而实践路径高度收缩”的现象\cite{dimaggio1983iron, meyer1977institutionalized, weick1976educational}；复杂系统理论则提醒我们，高校并非简单的机械系统，而是由多种异质主体构成的复杂适应系统，其宏观行为往往由微观互动的非线性涌现所塑造\cite{holland1995hidden, miller2007complex}。

然而，从研究现状看，仍存在三个明显不足：第一，在理论逻辑上，多数分析仍然采用\textbf{线性因果}视角，倾向于从单一维度（如经费、政策、结构）寻找“关键变量”，对多主体、多层级互动带来的非线性效应关注不足；第二，在方法上主要依赖问卷、访谈与统计回归，难以系统刻画教师、学院与职能部门在长期博弈中的策略演化与相互影响；第三，对“临界阈值”“路径锁定”“相变”等复杂系统特征缺乏情景化、可视化的分析工具。

随着大数据、复杂网络和智能体模拟等方法的发展，计算社会科学与计算教育学为重新审视教育治理问题提供了新的路径\cite{lazer2009life, epstein2006generative, zhang2023computational}。特别是近年来大语言模型及其驱动的智能体技术快速发展，使得基于文本、访谈和政策文件构建“类人行动者”的数字化表达成为可能，为在虚拟空间中搭建“人造高校治理系统”、探索不同制度参数下的长期演化路径提供了技术基础\cite{park2023generative, hong2023metagpt, sumers2023cognitive, wang2024survey}。

在此背景下，本文聚焦三个相互关联的核心问题：

\begin{enumerate}
    \item 在现有评价体系与科层结构下，地方高校应用型转型迟滞的\textbf{微观动力机制}是什么？教师、中层干部等行动者如何在有限理性和既有规范约束下进行策略选择？
    \item 转型改革中是否存在可识别的\textbf{临界阈值}与\textbf{相变机制}？例如，应用型成果在评价体系中的权重调整到何种程度，才可能引发教师行为从“重论文”向“重应用”的集体转向？
    \item 如何利用大语言模型驱动的智能体与数字孪生技术，在虚拟情境中对不同改革方案进行\textbf{反事实推演}，为现实中的政策设计提供可供参考的“计算实验室”？
\end{enumerate}

围绕上述问题，本文在国家教育科学规划课题（B1A250166）长期田野调研和数据积累基础上，构建嵌入大语言模型智能体的“高校治理数字孪生系统”，在复杂适应系统框架下，通过演化博弈与计算实验生成性地分析地方高校应用型转型的动力机制。下一章将从理论视角出发，系统梳理地方高校转型、组织制度与复杂系统等相关研究，并提出本文的综合分析框架。

% =========================================================
% 第二章 理论视域 (Theoretical Framework)
% =========================================================
\section{理论视域：理论丛林中的路径选择与重构}

\subsection{地方高校应用型转型的研究图景与不足}

围绕地方本科院校向应用型大学转型，国内外研究已形成较为清晰的宏观图景。潘懋元、钟秉林、瞿振元等学者从体系结构、类型定位与治理改革等维度指出：地方高校在提升高等教育毛入学率、支撑区域发展方面具有基础性地位，转型是缓解“多而不优”“大而不强”结构性矛盾的必由之路\cite{pan2014transformation, zha2011understanding}。在类型定位上，“应用型大学”被界定为区别于研究型大学和高职院校的独立类型，其人才培养目标由“学术再生产”转向“实践能力与职业胜任力”\cite{qu2014connotation}。从政策改革看，学者们也普遍注意到，在市场化、分权化与全球化的综合背景下，中国高等教育治理结构重塑给地方高校带来了新的压力与机遇\cite{mok2005globalization}。

另一条研究线索则聚焦于地方高校转型的\textbf{制度障碍}。Cai 等从制度变迁视角指出，地方高校在争取资源与合法性的过程中出现明显的“制度锁定”：一方面，“211”“985”“双一流”等精英导向项目在资源配置上强化了金字塔式结构；另一方面，以科研论文、学位点数量为中心的评价体系持续加固“向上模仿”的路径依赖，使得服务地方的功能被系统性弱化\cite{cai2014institutional}。国内研究也揭示了“破五唯”之前，地方高校科研评价中长期存在“唯论文”“唯项目”“唯帽子”等倾向，导致“重学术、轻应用”的激励扭曲\cite{wu2019evaluation}。

总体而言，既有研究在三个方面为本文奠定了重要基础：第一，围绕“地方本科院校转型发展”的宏观政策取向已较为清晰，为本文界定问题提供了规范性参照\cite{pan2014transformation, zhong2014applied}；第二，关于“应用型大学”的内涵、类型标准与建设路径已有较为系统的理论阐释，为理解“转向何处”提供了价值坐标\cite{qu2014connotation}；第三，制度分析与政策研究揭示了评价体系、资源分配与治理结构对转型的制约机制，为本文进一步追问“为何难转、如何能转”提供了结构性线索\cite{mok2005globalization, cai2014institutional}。

与此同时，从理论和方法上看，现有研究仍存在三方面不足：其一，多数研究停留在\textbf{静态结构分析}层面，着力描述政策文本、组织结构与指标体系，对转型过程中的\textbf{动态演化机制}涉及有限；其二，尽管新制度主义和组织社会学已经被引入，但更多用于解释“趋同化”“合法性追求”等宏观现象，对教师、院长等微观主体在科层制约束下的\textbf{认知博弈与策略选择}关注不足\cite{dimaggio1983iron, meyer1977institutionalized}；其三，在方法上主要依赖问卷调查、访谈与统计回归，难以应对地方高校转型这一\textbf{典型复杂适应系统}中的非线性效应、路径依赖与相变现象\cite{holland1995hidden, miller2007complex}。

随着计算社会科学和计算教育学的兴起，越来越多研究开始尝试使用基于智能体的仿真、复杂网络分析等方法研究教育系统\cite{lazer2009life, epstein2006generative, zhang2023computational}。但就公开文献而言，将大语言模型驱动的智能体与中国地方高校应用型转型情境深度结合的工作仍然稀缺。这恰恰为本文提供了切入空间：在既有宏观结构研究的基础上，引入\textbf{复杂适应系统 + 组织社会学 + 生成式计算实验}的综合视角，重构地方高校转型困境的理论图景。

\subsection{理论视角的整合：从资源、制度到复杂系统}

为理解地方高校为何长期徘徊于“形转实不转”的困境，有必要对相关理论传统进行有选择的整合，而非简单并列。

\subsubsection{资源依赖与环境压力}

资源依赖理论将组织视为深度嵌入资源环境的开放系统，组织行为在很大程度上是对关键资源控制者偏好的回应\cite{arthur1999complexity}。在地方高校情境中，教育行政部门、地方政府及用人单位构成了核心环境：前两者掌握经费拨付、学位点审批和职称评审权，后者则通过人才需求结构传递“应用型”信号。从这一视角看，地方高校倾向于“迎合资源配置者偏好”，在现实中更重视能够直接带来财政拨款、项目和头衔的指标（如 SCI 论文、纵向课题），而对横向项目、技术转移等应用型成果投入不足，成为可预期结果。

但资源依赖理论在解释转型困境时也存在局限：一是倾向于将组织视作“理性适应者”，假定其会根据环境变化做出收益最大化选择，低估了内部文化惯性与认知偏差；二是对多层级科层结构内部的权力博弈关注有限，难以刻画高校内部“校—院—系—个人”之间细腻的力量关系与信号传导。

\subsubsection{新制度主义与制度同形性}

新制度主义强调，组织的行为不仅是对资源约束的理性回应，更是对\textbf{制度环境中合法性期待}的象征性顺从\cite{dimaggio1983iron, meyer1977institutionalized}。在高教领域，制度同形性解释了地方高校为何不断模仿研究型大学：在强制度环境中，组织为获得合法性与声誉，会趋向采用被视为“先进”“正统”的结构与实践，即便这些安排并不契合自身使命。这一视角有助于理解地方高校在学科设置、科研评价、人才引进等方面的“千校一面”。

同时，Meyer 和 Rowan 提出的“形式结构的神话与仪式”提示，高校内部存在大量“脱耦合”现象：正式制度安排（如章程、办法）与日常实践之间往往存在显著差距，许多改革在文本上被贯彻，却在执行中被“仪式化”，只保留形式而未进入真实决策与行为层\cite{meyer1977institutionalized, weick1976educational}。在本文所关注的议题中，“转型写在规划里、论文写在考核里”正是这种脱耦合的典型呈现。

然而，新制度主义在处理\textbf{动态演化}时仍显不足：一方面，侧重揭示制度结构如何塑造行动者，对行动者如何通过策略性行为反过来重塑制度关注有限；另一方面，对非线性效应、临界阈值和多稳态等“复杂系统”特征涉及不多。

\subsubsection{复杂适应系统视角}

复杂适应系统（Complex Adaptive Systems, CAS）理论则将组织与社会视为大量相互作用的异质主体构成的开放系统，这些主体在局部信息和有限理性约束下做出决策，通过持续互动涌现出宏观模式\cite{holland1995hidden, miller2007complex}。高校具备 CAS 的典型特征：教师、学院、职能部门和上级主管单位在不同层级上进行策略选择，其行为同时受到制度规则、社会规范与个体认知的共同影响。

从 CAS 视角出发，地方高校转型困境不只是“资源不足”或“制度不佳”的静态结果，更是多主体在既有制度与文化背景下长期博弈所形成的\textbf{路径依赖与锁定效应}。例如，当“论文数量”长期被视为职称晋升的唯一“硬通货”时，即便政策逐步强调应用型成果，个体仍可能基于风险与不确定性考量而选择“路径保守”，系统也就难以自动运转至新均衡状态。这也意味着，转型改革很可能存在某种“临界阈值”与“相变”现象：在关键参数未跨越阈值前，系统表现出强惯性；一旦跨越，则可能出现行为的雪崩式变化\cite{scheffer2009early}。

因此，将资源依赖、新制度主义与 CAS 视角结合，可以形成一种更为立体的解释框架：外部资源与合法性压力提供约束边界，制度同形性与脱耦合塑造“形式结构”，而多主体在其中的互动与学习则决定系统最终落入何种稳定状态。

\subsection{“生成式治理”的概念界定与分析框架}

\subsubsection{从控制论治理到生成式治理}

传统教育治理在逻辑上更多建立在控制论式的“反馈—调节”模型之上：通过设定目标、设计方案、实施政策与收集反馈，构成“输入—处理—输出—反馈”的线性闭环\cite{wiener1948cybernetics}。这种模式在环境相对稳定、结构相对简单的情境下具有较强解释力，但在高度复杂、快速变化的教育系统中，其局限日益凸显：一是反馈往往存在信息滞后与信号失真，治理主体难以及时捕捉基层行动者的真实反应；二是线性调节机制难以应对非线性效应和相变现象，容易出现“政策一刀切”“激励失衡”等后果。

随着治理现代化与教育数字化进程推进，国内学者提出了“多元共治”“教育善治”“教育治理体系和治理能力现代化”等概念，强调通过多主体参与、程序正义与信息公开提升治理效能\cite{shi2020governance, chu2016governance}。与此同时，以 Lazer 等人为代表的计算社会科学研究提出，应借助大数据、仿真与可视化手段构建“人造社会”，在虚拟空间中对政策方案进行预演与评估\cite{lazer2009life, epstein2006generative}。教育领域也开始出现“计算教育学”的讨论，尝试用数据、模型和算法重构对教育现象的理解与干预路径\cite{zhang2023computational}。

在这一背景下，本文提出\textbf{“生成式治理”（Generative Governance）}概念，指的是：在高度数智化环境中，治理主体不再仅依赖经验判断和线性反馈，而是通过构建基于真实数据和智能体的\textbf{数字孪生系统}，在虚拟空间中重构微观主体的认知结构与行为规则，利用计算实验生成多种可能的制度配置及其长期后果，从而实现对现实改革的\textbf{预演、择优与精准实施}。如果说传统治理更像“事后校准的温控器”，生成式治理则更接近“事前调参的风洞实验”。

\subsubsection{“认知—制度—演化”综合分析框架}

在具体分析框架上，本文在资源依赖、新制度主义与 CAS 的基础上，构建了一个“认知—制度—演化”的综合图式，用以刻画地方高校应用型转型困境的生成逻辑：

\begin{enumerate}
    \item \textbf{认知层（Cognition）}：以教师、院长、中层干部等行动者为核心，关注其价值取向、风险偏好与有限理性\cite{simon1996sciences}。例如，“躺平型”教师更关注工作量与稳定性，“焦虑型”青年教师高度敏感于考核指标与晋升机会，“守门人型”中层干部在“维稳”与“改革”之间寻求平衡。这些差异化认知通过相互观察与模仿，会在组织内部形成特定的规范预期与参照框架。
    \item \textbf{制度层（Institution）}：包括职称评审办法、科研评价细则、绩效分配规则等正式制度，也包括“论文至上”“帽子崇拜”等非正式规范\cite{weber1978economy, meyer1977institutionalized}。制度层一方面为行动者提供激励与约束，另一方面也被行动者通过策略性行为不断再生产和改写，产生“形式结构的仪式化”和“执行层面的脱耦合”。
    \item \textbf{演化层（Evolution）}：关注在认知与制度互动下，行动者策略组合如何随时间演变，并最终涌现为宏观层面的系统状态\cite{holland1995hidden, miller2007complex}。在 CAS 视角下，这一过程往往呈现非线性特征，可能出现“激励倒 U”“行为锁定”“临界相变”等现象，其结果并非简单的线性叠加。
\end{enumerate}

在方法论上，本文通过构建嵌入大语言模型智能体的“高校治理数字孪生系统”，将真实调研中提炼的认知模式、制度规则与互动情景编码进模型，使上述三层结构在虚拟环境中得以“具象化”。在这一意义上，\textbf{生成式治理}并非仅是规范性口号，而是一套可计算、可试验、可迭代的治理实践框架：通过在数字孪生中持续进行反事实推演，观察不同政策参数下“认知—制度—演化”的联动效应，从而为现实中的地方高校应用型转型提供更加科学的政策依据。

与既有研究相比，本文的理论贡献主要在于：一是在地方高校应用型转型议题上引入 CAS 视角，将“转型难题”从静态结构问题重构为动态演化问题；二是提出并操作性界定“生成式治理”，将其与数字孪生和计算实验紧密耦合；三是通过大语言模型智能体建模，将教师认知、制度规则与组织演化整合到统一的计算框架之中，为教育治理研究提供了一种可复制的方法路径。基于此，下一章将详细介绍本研究的计算实验设计，包括数据基础、智能体架构与历史回溯检验等内容。



% =========================================================
% 第三章 研究设计与方法 (Research Design and Methods)
% =========================================================
\section{研究设计与方法：高校治理数字孪生与智能体建模}

\subsection{数据基础与情景抽取}

本研究的数据基础主要来自三类材料：一是依托全国教育科学规划课题（B1A250166）连续 10 年形成的纵向面板数据，涵盖某省若干地方本科高校在经费结构、科研产出、人才培养与社会服务等方面的关键指标；二是面向校领导、中层干部与一线教师开展的深度访谈与焦点小组，共计千余份有效文本；三是高校内部的规划文本、职称评审办法、绩效分配细则等制度性文件。

在数据预处理阶段，研究团队遵循“量化结构—情景抽取—行为画像”的思路，将多源材料转化为可建模的结构化信息：

\begin{itemize}
    \item 在量化结构层面，对面板数据进行清洗、缺失值处理与标准化，构建涵盖经费来源与用途、纵横向项目、论文与专利产出、横向合作等维度的指标体系，并通过主成分分析与聚类分析识别不同发展路径的高校类型。
    \item 在情景抽取层面，结合文本编码与扎根理论，对规划文本、制度文件和访谈资料进行开放式编码、主轴编码与选择性编码，提炼出“职称评审中论文与应用成果权重设定”“横向课题折算系数”“中层干部对改革文件的再解释策略”等高频情景单元。
    \item 在行为画像层面，围绕教师、中层干部和校领导三个关键行动者群体，抽取其价值取向、风险偏好、行动约束与策略倾向，形成一组具有代表性的“认知原型”，如“论文导向型青年教师”“项目驱动型骨干教师”“维稳优先型中层干部”等。
\end{itemize}

上述工作为后续智能体建模提供了真实语境与边界条件，使数字孪生系统在建构过程中尽可能保持对高校治理情景的“语境忠实”。

\subsection{智能体类型与认知建模}

在智能体建模层面，本文采用“角色原型 + 大语言模型 + 规则约束”的混合范式，将高校治理中的主要行动者抽象为若干类智能体，并通过大语言模型重构其认知与决策过程\cite{park2023generative, hong2023metagpt, sumers2023cognitive, wang2024survey}。

\subsubsection{教师智能体}

教师智能体是系统的最基本单元，其状态变量主要包括：学科类别、职称层级、年龄段、科研基础、家庭负担、职业期望等；行为变量则包括：时间分配（教学、科研、社会服务）、项目偏好（纵向/横向）、成果类型选择（论文、专利、报告、技术转移）、校内外合作意愿等。

在认知建模上，本文基于大语言模型构建“教师认知子模型”。该子模型的输入包括：学校当前制度参数（如职称评审权重、绩效分配规则）、同行行为分布、上级文件中的政策信号等；输出则为教师在给定情景下的“主观理解”与“策略选择”，例如：“将新出台的横向课题加分政策视为短期口号，维持既有论文优先策略”，或“意识到应用型成果在晋升中的实际权重提升，主动寻找企业合作机会”。

为避免自然语言生成导致的过度发散，模型在生成教师行为时叠加一组参数化的“行为约束规则”：包括总工作时间上限、不同类型成果的最低时间成本、重大策略调整的“心理摩擦”与“信息滞后”等。这些规则与大语言模型输出共同决定教师智能体在每一考核周期内的具体行动。

\subsubsection{中层干部智能体}

中层干部智能体主要包括学院院长、系主任和职能部门负责人等，是连接学校决策与一线教师的关键“中介层”。一方面，他们需要回应学校和上级主管部门的考核要求；另一方面，又要维护学院内部的运转稳定与可接受度，具有典型的“夹心层”特征。

在建模中，中层干部智能体的状态变量包括：所在学院资源禀赋、学生规模、历史绩效、人员结构及个人晋升路径等；行为变量包括：如何解读与转述上级改革文件、如何制定学院内部实施细则、如何在执行过程中进行选择性强调或弱化等。大语言模型在此扮演“政策翻译器”的角色：同一份学校层面的改革文件，在不同中层干部智能体中可能被解读为“实质性机会”或“形式性任务”，进而造成截然不同的执行路径。

\subsubsection{学校管理者智能体}

学校管理者智能体主要指校领导层，其决策集中体现在学校层面的发展规划、资源配置和制度设计上。相较于教师与中层干部，学校管理者更关注外部环境（如地方政府、上级部门、竞争高校）的评价与期待，更敏感于学校整体排名、声誉与“综合办学绩效”。

在模型中，学校管理者智能体通过调整若干关键制度参数来影响系统运行，例如：职称评审中应用型成果折算系数、绩效分配中横向经费奖励倍数、对产业学院建设的资源倾斜力度等。其决策同样由大语言模型驱动，在输入中综合考虑上级政策文件、外部排名压力和内部反馈信息，并在规则约束下形成可执行的制度调整方案。

\subsection{策略空间与激励机制编码}

为开展系统性实验，本研究在制度层将复杂多样的政策工具抽象为若干可调节的“激励参数”，主要包括：

\begin{itemize}
    \item 职称评审中不同成果类型的等效权重（如论文、专著、技术报告、专利、横向课题、社会服务等）。
    \item 绩效分配中纵向经费与横向经费的奖励倍数以及基础绩效占比。
    \item 教师工作量考核中教学、科研与社会服务三类任务的最低要求与上限配置。
    \item 对产业学院、校企合作平台、实践基地等的资源倾斜强度与配套政策。
\end{itemize}

在策略空间层面，教师与中层干部智能体可在给定激励参数下选择不同策略组合。例如，教师可以在“论文优先”“论文—应用平衡”“应用倾斜”等策略之间切换；中层干部可以在“刚性执行”“选择性执行”“符号执行”等模式之间选择。每种策略组合在一个考核周期内将产生不同的系统输出，包括论文与项目数量、横向经费占比、学生实践学分比例等。

为避免策略空间过于庞大导致状态爆炸，本文采用“原型策略 + 局部微调”的方式：先根据实地调研中提炼的典型行为模式构建有限个基础策略，再允许智能体在局部范围内进行连续调整，从而在可控复杂度下保持足够的行为多样性。

\subsection{高校治理数字孪生环境与情景设定}

在上述智能体与策略空间基础上，本文构建了一个多层级的“高校治理数字孪生环境”，并在其中设定若干具有代表性的情景，用于开展计算实验。

\subsubsection{环境结构与时间步长}

数字孪生环境包含三个层级：学校层、中层学院层与教师个体层。时间以“考核周期”为基本步长，每个周期对应 1 年。每一时间步系统更新流程如下：

\begin{enumerate}
    \item 学校管理者智能体根据外部环境和内部绩效反馈调整激励参数；
    \item 中层干部智能体对新的制度参数进行解读，形成学院层实施细则与操作指引；
    \item 教师智能体在感知到的激励结构与同行行为下更新策略选择，决定时间分配与项目偏好；
    \item 系统汇总各类成果产出与行为数据，形成该周期的“绩效向量”，并作为下一轮决策的输入。
\end{enumerate}

其中，第 2 步和第 3 步的决策过程由大语言模型生成的自然语言“意向”与预设规则共同决定，再映射为结构化的策略变量和产出指标，从而在保证行为可解释性的同时维持模型的可计算性。

\subsubsection{情景设定与政策变量}

为检验不同政策组合对系统演化路径的影响，本文设置若干典型情景，包括但不限于：

\begin{itemize}
    \item 现状情景：保持当前职称权重与绩效规则不变，用于校准与验证模型对现实系统的拟合能力。
    \item 温和改革情景：适度提高应用型成果权重与横向经费奖励倍数，但中层干部考核机制基本维持不变。
    \item 深度联动情景：同步调整教师职称评价、学院绩效考核与学校外部评价，将应用型成果权重提升到与论文相当的区间，并嵌入中层干部激励约束。
    \item 过度激励情景：显著提高横向经费与技术转移奖励倍数，以检验是否会触发古德哈特效应和短期逐利行为。
\end{itemize}

在每个情景下，系统运行多个时间步，并通过多次蒙特卡洛重复实验观测行为分布、绩效指标与组织状态的演化轨迹，从中识别稳定模式、临界阈值与相变现象。

\subsection{技术路线、关键技术和主要创新点}

综合上述设计，本研究的技术路线可概括为“实证调研—认知抽象—智能体构建—数字孪生—计算实验—政策反演”六个环节。围绕这一路线，关键技术与主要创新点集中体现在以下三个方面：

\begin{itemize}
    \item 在认知建模上，将地方高校教师与中层干部的“主观世界”嵌入大语言模型，通过提示工程、示例对话与规则约束构建具备情境感知与角色一致性的“认知智能体”，较好缓解了传统仿真模型中对行动者过度简化的问题。
    \item 在系统构建上，提出面向高校治理的数字孪生框架，将制度参数、行为策略与绩效输出映射到统一状态空间，使复杂治理过程实现可视化、可计算与可复现，为政策情景的批量推演提供“虚拟场域”。
    \item 在政策分析上，首次尝试以生成式计算实验的方式识别地方高校应用型转型中的临界阈值与相变机制，并据此提出“非线性精准激励”思路，将传统经验式决策升级为基于模型证据的“风洞试验式”决策。
\end{itemize}

\subsection{技术路线可行性与先进性分析}

从可行性看，一方面，本研究依托长期积累的面板数据与扎实田野调研，为认知抽象与参数设定提供了可靠基础；另一方面，大语言模型与多智能体研究在生成逼真对话、重构社会互动等方面已取得显著进展，为本研究提供了成熟技术工具\cite{park2023generative, hong2023metagpt, sumers2023cognitive}。此外，研究团队兼具教育治理研究与计算建模经验，可支撑数字孪生系统的迭代开发与情景扩展。

从先进性看，本研究突破了传统教育政策研究“描述—解释—建议”的线性范式，将地方高校转型置于 CAS 框架下，通过“人造高校治理系统”探索制度变革的长期后果。在方法论上，借助大语言模型对行动者认知进行显性建模，使以往难以量化的“制度理解偏差”“风险感知”“同行效应”等因素得以进入计算分析。在理论上，以“生成式治理”概念串联资源依赖、新制度主义与复杂系统理论，为理解和推动教育治理现代化提供了新的分析视角与方法支撑。

% =========================================================
% 第四章 计算实验结果与生成式发现 (Findings)
% =========================================================
\section{计算实验结果与生成式发现}

\subsection{基线情景：现状复刻与模型校准}

在保持现有职称权重与绩效规则不变的基线情景下，数字孪生系统的演化结果与实证数据在若干关键维度上高度一致：第一，论文数量与纵向课题经费保持温和增长，而横向经费占比和技术转移指标变化缓慢；第二，教师时间分配中，科研尤其是论文写作占据主导位置，参与企业合作和社会服务的比例偏低；第三，中层干部在传达改革文件时普遍存在“强调可量化指标、弱化制度机制”的倾向，更关注短期可考核结果。

这一情景下，系统表现出明显的路径依赖特征：即便外部环境出现适度变化（如地方政府增加对应用型成果的宣传和奖励），内部行为模式仍保持较强惯性，难以自动发生根本改变。这一结果表明，模型在刻画“以论文为核心”的行为均衡方面具有较好拟合能力，也为后续情景分析提供了参照基准。

\subsection{激励强度与古德哈特效应的倒 U 关系}

在逐步提高横向经费和应用型成果奖励倍数的情景下，系统呈现出典型的“倒 U 型”性能曲线：当激励强度从较低水平提升到中等区间时，教师参与企业合作和技术转移的积极性显著提高，相关指标随之上升；但当激励强度持续增加超过某一节点后，系统开始出现古德哈特效应——“当一个指标成为目标，它就不再是一个好的指标”。

具体而言，部分教师与中层干部智能体发展出针对指标的策略性行为，例如优先选择风险较低、技术含量有限但容易形成合同金额的横向项目，或将原本应通过纵向课题完成的工作拆解为多个形式化横向合作，以最大化绩效得分。长期运行后，应用型成果数量虽然增加，但质量与对地方产业的真实贡献并未同步提升，甚至出现“虚高”的风险。

图 \ref{fig:goodhart} 给出了激励强度与应用型成果有效贡献之间关系的示意图。可以看到，在过低激励区间，教师缺乏足够行动动机；在适中激励区间，数量与质量相对协调；而在过高激励区间，成果数量继续上升但有效贡献开始下滑。

\begin{figure}[htbp]
    \centering
    \fbox{\parbox{0.8\textwidth}{\centering \vspace{2cm} \textbf{[图2]：激励强度与应用型成果有效贡献的倒 U 关系示意}\\ \small （示意图：横轴为激励强度，纵轴为有效贡献，曲线呈倒 U 型） \vspace{2cm}}}
    \caption{激励强度与古德哈特效应的倒 U 关系}
    \label{fig:goodhart}
    \vspace{-0.5em}
    {\footnotesize \kaishu 注：有效贡献综合考虑技术含量、转化率与地方满意度等因素。}
\end{figure}

这一结果提示，应用型转型中的激励设计并非“越强越好”，而需要在数量与质量、短期与长期之间寻求平衡，通过多维指标组合与过程性评价降低古德哈特效应风险，使激励强度保持在“倒 U 曲线”的合理区间。

\subsection{评价权重阈值与教师行为相变}

在逐步提高应用型成果在职称评审中的等效权重的情景下，系统呈现出清晰的“临界阈值”特征：当应用型成果权重处于较低水平时，大多数教师仍维持以论文为主的策略，仅少数具有企业合作优势的教师增加应用型投入；当权重提升至约 30\%--40\% 区间时，教师群体的策略选择出现集体性“相变”，大量教师开始重新分配时间，在保证基本论文产出的前提下，将更多精力投入横向课题和技术服务。

这一相变并非线性叠加的结果，而是多重机制共同作用的体现：一方面，在“论文至上”的规范仍占主导时，早期转向应用型的教师面临较高声誉风险和晋升不确定性，难以形成可模仿的示范效应；另一方面，当权重跨过关键阈值后，教师通过日常交流逐渐形成对“制度变迁真实程度”的共识，对未来收益结构的预期发生改变，从而在较短时间内形成新的集体行为规范。

图 \ref{fig:phase-transition} 展示了不同权重水平下教师策略分布演化趋势的示意。可以看到，在临界区间附近，应用型策略比例出现明显“拐点”，与复杂系统中的相变现象具有高度相似性。

\begin{figure}[htbp]
    \centering
    \fbox{\parbox{0.8\textwidth}{\centering \vspace{2cm} \textbf{[图3]：职称评价权重调整与教师策略分布相变示意}\\ \small （示意图：横轴为时间步，纵轴为不同策略教师比例，在临界权重区间附近出现突变） \vspace{2cm}}}
    \caption{职称评价权重调整与教师行为相变}
    \label{fig:phase-transition}
    \vspace{-0.5em}
    {\footnotesize \kaishu 注：不同颜色代表“论文优先”“平衡发展”“应用倾斜”等策略类型。}
\end{figure}

这一发现表明，地方高校应用型转型不仅需要“释放改革信号”，更需要通过权重调整跨越关键阈值，促成教师群体预期结构与行为模式的整体重构。

\subsection{中层过滤、信息熵增与“最后一公里”}

在多个情景实验中，中层干部智能体的行为被证明是影响改革落地效果的关键变量。即便学校层面通过调整职称评价与绩效分配实现了制度参数的实质变动，如果中层干部在解读与传达过程中出于风险规避考虑，将改革内容“稀释”为“暂可观望”“先看执行效果再说”，则一线教师对制度变化的感知会显著弱化。

模型显示，在未同步调整中层干部考核机制的情景下，改革信号在传导过程中会经历明显的信息熵增：原本清晰的改革目标在层层传递中逐渐模糊，甚至被其他短期任务淹没。其结果是：学校层面的制度调整“存在于文件中”，但教师层面的行为模式“停留在旧均衡”，形成典型的“文件改革”而非“行为改革”。

相反，在将中层干部考核指标与学院应用型成果直接绑定，并通过数字孪生模拟向其展示不同执行策略长期后果的情景下，中层干部更倾向于采取实质性执行路径，主动设计符合学院特点的应用型转型方案。此时，信息在传导过程中的熵增显著降低，改革信号得以较为完整地传递至一线教师。

这一结果进一步印证了“生成式治理”强调的一个核心判断：在高度科层化的教育系统中，单向命令—控制难以奏效，必须通过对关键中介节点（尤其是中层）的激励重构与认知重构，打通政策执行的“最后一公里”。

\subsection{生成式政策启示的初步提炼}

综合上述计算实验，可以初步提炼出若干具有操作性的政策启示：

\begin{itemize}
    \item 在激励设计上，应避免单一指标的过度放大，通过多维度组合指标与过程性评价降低古德哈特效应风险，使激励强度保持在“倒 U 曲线”的合理区间。
    \item 在权重调整上，应有意识地跨越教师群体行为相变的关键阈值，使改革不再停留于“边际微调”，而是通过“可识别的转折点”重塑集体预期。
    \item 在治理结构上，应将中层干部视为“放大器”而非“过滤器”，通过考核绑定、能力支持与可视化沙盘推演，促使其从“风险规避型守门人”转变为“应用型转型的共创者与设计者”。
    \item 在决策方式上，应常态化运用数字孪生与计算实验，将政策设计置于“可视—可算—可迭代”的政策风洞之中，以生成式方式持续优化，而非一次性“定型”。
\end{itemize}

这些启示将在下一章中与案例高校的实证材料相互对话，进一步凝练为面向地方高校的治理策略组合。

% =========================================================
% 第五章 实证对话与治理策略生成
% =========================================================
\section{实证对话与治理策略生成}

\subsection{案例高校对照与模型校准}

为避免数字孪生系统沦为“纯理论演绎”，本研究选取若干具有代表性的地方高校开展案例对照，将计算实验结果与实地调研材料进行互证与校准。

\subsubsection{类型案例的选择逻辑}

在案例选择上，重点考虑以下维度的差异性与代表性：

\begin{itemize}
    \item 区域属性差异：既包括沿海发达地区的地方高校，也包括中西部欠发达地区的地方高校，以观察不同区域发展环境下转型路径的差异。
    \item 办学基础差异：既包括原本科院校，也包括由高职专科升格或通过合并重组形成的地方高校，以比较不同历史起点的约束条件。
    \item 转型阶段差异：既包括已获批省级或国家级应用型示范高校的“先行者”，也包括刚刚启动转型规划的“追赶者”，以识别不同阶段的治理重心。
\end{itemize}

在此基础上，构建若干类型化案例组合，如“区域资源充沛但内部治理惯性较强型”“区域环境一般但校内改革动力充沛型”等，为数字孪生情景设定提供现实参照。

\subsubsection{实地调研与计算结果的映照}

在每所案例高校中，研究团队重点围绕以下问题开展访谈与文本分析：

\begin{itemize}
    \item 学校近年来在职称评价、绩效分配、产业学院建设等方面的具体改革举措及时间脉络。
    \item 教师和中层干部对相关改革的主观感受及行为调整情况，包括“是否真切感受到应用型成果价值提升”“是否认为横向项目‘值得投入’”等。
    \item 改革过程中出现的预期外后果，如形式主义横向项目、教学任务被过度边缘化、管理负担显著增加等。
\end{itemize}

随后，将案例高校的制度参数与行为特征输入数字孪生系统，运行对应情景的计算实验，对比模型输出的绩效向量与实测数据的差异。例如，当模型预测在某种激励结构下应出现显著行为相变，而实地并未观察到时，提示可能存在模型尚未充分刻画的约束因素（如地方政府高压考核、学科评估硬性指标等），需要在后续建模中予以补充。

通过这一往返循环，数字孪生系统不再只是抽象推演工具，而逐步成为与现实高校治理实践持续对话的“第三空间”。

\subsection{生成式策略菜单的构建}

在充分比对计算实验与案例材料的基础上，本研究不试图给出单一“最佳方案”，而是构建面向不同类型高校的“策略菜单”，供决策者依据自身情境进行组合选择与渐进实施。

\subsubsection{职称评价与绩效分配的组合路径}

围绕职称评价权重与绩效分配规则的联动设计，可以抽象出若干具有代表性的策略组合路径，例如：

\begin{itemize}
    \item 渐进调适路径：在保持论文基本权重的前提下，分阶段提高应用型成果等效分值，并设置过渡期“成绩豁免”或“缓冲期”机制，使教师在尝试新路径时避免遭遇即时惩罚。
    \item 阈值跨越路径：在充分沟通与预告基础上，在某一年度将应用型成果权重一次性提升至临界区间，同时配套完善评审标准与典型案例，用“清晰转折点”取代“长期模糊期望”。
    \item 差异化路径：针对不同学科、不同发展阶段教师群体设置差异化权重结构，例如对工科与产业关联度高的专业提高横向项目和技术报告权重，对基础学科更多强调高质量论文与高水平科研项目。
\end{itemize}

数字孪生系统可以预先模拟上述不同路径在多个考核周期内的绩效演化、教师行为变化及潜在风险点，为高校决策层提供“制度参数—行为反应—结果后果”之间的可视化映射。

\subsubsection{中层干部激励与支持机制的设计}

针对中层干部在改革中的关键角色，策略菜单主要包括：

\begin{itemize}
    \item 目标绑定机制：将学院层应用型成果指标与中层干部年度绩效和任期考核紧密挂钩，同时避免单一指标，通过多维度组合减少投机空间。
    \item 能力支持机制：提供面向中层干部的“应用型转型工作坊”“数字孪生沙盘演练”等平台，使其在安全环境中进行策略试错，提升制度设计与沟通能力。
    \item 反馈迭代机制：建立自下而上的反馈渠道，将一线教师对改革措施的体验与改进建议纳入学院与学校层面的决策循环，削弱“单向命令”带来的僵化效应。
\end{itemize}

通过在数字孪生环境中引入不同执行策略的模拟，可以直观呈现“符号执行”“选择性执行”与“实质执行”在长期后果上的显著差异，帮助校领导层认识到“激励不到中层，改革难落地”的结构性问题。

\subsection{推广条件、边界与情境敏感性}

本研究提出的治理策略强调情境敏感性与边界意识，并不主张“一刀切”的普适方案。

\subsubsection{推广的必要前提}

从计算实验与案例对比来看，策略菜单的可行推广至少需要满足以下前提条件：

\begin{itemize}
    \item 高校在数据治理层面具备基本条件，能够持续、规范地收集与整理教学、科研和社会服务等关键指标，为数字孪生模型更新提供支撑。
    \item 高校领导层对应用型转型具有相对稳定的政策意志，愿意接受“短期指标波动”以换取“中长期结构优化”，避免在改革初期因个别指标下降而仓促“急刹车”。
    \item 教师群体具备一定的参与空间，能够通过民意调查、教学委员会、学术委员会等渠道对新制度提出反馈，避免将改革理解为单纯“自上而下的任务增加”。
\end{itemize}

\subsubsection{策略应用的风险与边界}

同时，策略菜单也存在明确的适用边界。例如，在某些学科供给严重过剩、就业吸引力明显不足的情境下，即便调整职称权重和绩效规则，也难以仅依靠内部激励实现高层次应用型成果突破；在地方产业结构单一、企业合作基础薄弱的地区，过度强调横向经费可能诱发“虚假合作”与“指标空转”。

因此，本文强调：数字孪生与计算实验应被视作辅助决策工具，而非替代决策的“自动化装置”。教育治理始终需要保留关于公共价值、长期使命与公平正义的规范性判断空间。

% =========================================================
% 第六章 结论、局限与未来展望
% =========================================================
\section{结论、局限与未来展望}

\subsection{研究结论的综合梳理}

围绕“地方高校应用型转型困境的形成机制及其治理策略”这一核心问题，本文在前期实证研究基础上，引入大语言模型与多智能体仿真，构建高校治理数字孪生系统，并通过一系列计算实验与案例对照，形成以下主要结论：

\begin{itemize}
    \item 应用型转型困境并非单一因素所致，而是职称评价、绩效分配、外部考核与内部认知长期耦合的结果，表现为制度激励与教师日常行为之间存在显著“错位”。
    \item 在激励设计方面，存在明显古德哈特效应风险：当单一指标被过度放大时，教师与中层干部容易发展出针对指标的策略性行为，导致数量上升而质量不升反降。
    \item 在职称评价权重调整过程中，教师群体行为转向呈现“相变”特征：只有当应用型成果权重跨越某一关键阈值后，集体预期才会发生实质改变，进而触发系统性行为重构。
    \item 中层干部在改革传导中发挥关键“过滤”或“放大”作用；若缺乏相应激励与能力支持，改革信号易在层层传递中发生信息熵增，导致“文件改革”难以转化为“行为改革”。
    \item 数字孪生与计算实验为教育治理提供了一种新的工具理性，使高校能够在政策实施前，在虚拟环境中预演不同制度组合的长期后果，以生成式方式不断修正与优化治理结构。
\end{itemize}

\subsection{理论与实践贡献}

在理论层面，本文尝试将资源依赖、新制度主义与复杂适应系统理论整合到“生成式治理”框架中，将教育治理理解为多主体在有限理性与互动网络中的持续生成过程，而非静态结构调整。通过引入大语言模型对行动者认知进行建模，本文为解释“制度—行为—绩效”之间的非线性关系提供了新的分析路径。

在实践层面，本文为地方高校提供了一套可直接嵌入校本实践的分析与决策支持工具：通过构建高校治理数字孪生模型，高校可以以较低边际成本模拟不同改革方案的中长期效果，识别潜在风险与临界阈值，从而在改革前期进行必要预沟通与配套设计，提升政策的可行性与稳健性。

\subsection{研究局限与反思}

尽管本研究在方法与框架上进行了探索性创新，但仍存在若干局限，需要在后续研究中加以弥补：

\begin{itemize}
    \item 在数据层面，尽管整合了多源材料，部分关键变量（如教师真实工作时间分配、项目深度参与程度等）仍依赖抽样调查与主观估计，可能导致模型参数存在偏差。
    \item 在建模层面，大语言模型对行动者认知的重构仍基于文本与有限样本，难以完全涵盖真实行动者复杂的心理与情感过程，存在“理性简化”的倾向。
    \item 在推广层面，数字孪生系统目前主要在有限数量案例高校中进行试验，其在更大范围内推广应用的条件与障碍尚需进一步检验与总结。
\end{itemize}

这些局限提醒我们：任何形式化模型都只能是对现实世界的“有损压缩”，其价值不在于给出唯一答案，而在于提供新的观察视角与对话平台。

\subsection{未来研究方向与工作展望}

围绕本文提出的问题与方法，未来研究可以从以下方向进一步拓展：

\begin{itemize}
    \item 在数据层面，逐步构建覆盖更多省份、更多类型地方高校的长时段面板数据库，引入学生学习成效、毕业去向质量、地方产业反馈等指标，以更全面评估应用型转型成效。
    \item 在方法层面，进一步发展“人机共建”的智能体建模范式，将教师与中层干部作为“共创者”纳入模型校准过程，通过工作坊形式共同修订行为规则与情景设定，提升模型解释力与认同感。
    \item 在工具层面，开发面向高校管理者的可视化决策支持平台，将复杂计算结果转化为直观图表与情景剧本，降低使用门槛，使数字孪生从研究工具演变为日常治理工具。
    \item 在比较研究层面，将地方高校应用型转型的治理经验与职业教育、高职本科以及国际类似院校进行系统比较，提炼具有跨情境参考价值的“生成式治理”一般原则。
\end{itemize}

总体而言，本文只是一个起点。地方高校应用型转型作为教育现代化的重要组成部分，将长期处于动态演化之中。如何在不确定的外部环境与多元诉求的内部结构之间寻求新的平衡，有赖于研究者与实践者在“生成式治理”理念下持续对话与共创。



% =========================================================
% English Title, Abstract and Keywords (for journal requirement)
% =========================================================

\clearpage
\begin{center}
\zihao{4}\bfseries
Generative Governance: A Complex Systems Analysis of Application-Oriented Transformation Dilemmas in Local Universities\\[0.5em]
\zihao{5}A Digital Twin–Based Computational Experiment with Large Language Model Agents
\end{center}

\vspace{0.8em}

\noindent\zihao{5}\textbf{Abstract}\quad
Guiding local universities to shift toward application-oriented development is a key link in building an education powerhouse and serving high-quality regional development. In practice, however, a long-standing paradox persists: strategic plans repeatedly stress “serving local development” and “industry–education integration”, while core mechanisms such as academic promotion and research evaluation remain dominated by papers and academic titles. As a result, transformation often becomes symbolic rather than substantive. Existing studies, which mainly adopt linear, policy- or structure-based causal analyses, struggle to reveal the nonlinear interaction between micro-level cognitive games and macro-level institutional rules when local universities are viewed as complex adaptive systems.

Drawing on ten-year longitudinal panel data and more than one thousand in-depth interviews from a national education research project (B1A250166), this paper adopts a complex adaptive systems and computational education science perspective to construct a “university governance digital twin” embedded with large-language-model-based agents. In a virtual environment, we conduct counterfactual policy simulations around typical reform scenarios. The computational experiments show that: first, simply amplifying the unit reward of academic performance easily triggers a Goodhart-type effect, with an inverted-U relationship between incentive strength and substantive innovation output; second, when the equivalent weight of application-oriented achievements in promotion criteria rises to roughly 30\%–40\% and receives explicit endorsement from middle managers, faculty behaviour undergoes a collective phase transition from “publication-oriented” to “application-oriented”, revealing an identifiable critical threshold; and third, along hierarchical transmission chains, middle managers often filter and reframe reform signals out of risk aversion, leading to a marked attenuation of policy intensity and information “entropy increase” in the “last mile” of implementation.

On this basis, the paper proposes a shift toward \emph{generative governance}: moving from experience-based judgement to a “policy wind tunnel” grounded in digital twins and computational experiments. Through non-linear, targeted incentives, data penetration across organisational layers, and the binding of middle-management interests to transformation outcomes, we reconstruct the motivational structure of application-oriented transformation and offer a computable, visualisable, and iterable decision-support framework for the modernisation of the education governance system and its governance capacity.

\vspace{0.8em}
\noindent\zihao{5}\textbf{Keywords}\quad
generative governance; application-oriented transformation; local universities; complex adaptive systems; computational education science; modernisation of education governance; digital twin; large language model agents


% =========================================================
% 参考文献 (Bibliography)
% =========================================================
\vspace{1cm}
\begin{spacing}{1.2} % 参考文献行距稍紧
\zihao{5} % 参考文献通常使用五号字
\bibliography{refs}
\end{spacing}


\end{document}